\documentclass[11pt]{report}
\usepackage[utf8]{inputenc}
\usepackage[T1]{fontenc}
\usepackage{lmodern}
\usepackage[english]{babel}
\usepackage{graphicx}
\usepackage{hyperref}
\usepackage{amsmath}
\usepackage{xcolor}
\usepackage{listings}
\usepackage{longtable}
\usepackage{booktabs}
\usepackage{geometry}
\geometry{margin=0.85in}
\setlength{\parskip}{2pt}

% --- Code listing style ---------------------------------------------------
\definecolor{codebg}{rgb}{0.96,0.96,0.96}
\definecolor{codecomment}{rgb}{0.4,0.45,0.4}
\definecolor{codestring}{rgb}{0.5,0.1,0.1}
\lstdefinestyle{shell}{
    backgroundcolor=\color{codebg},
    basicstyle=\ttfamily\small,
    breaklines=true,
    frame=single,
    numbers=none,
    language=bash,
    commentstyle=\color{codecomment},
    keywordstyle=\bfseries,
    aboveskip=6pt,
    belowskip=6pt
}
\lstdefinestyle{code}{
    backgroundcolor=\color{codebg},
    basicstyle=\ttfamily\small,
    breaklines=true,
    frame=single,
    numbers=left,
    numberstyle=\tiny\color{gray},
    commentstyle=\color{codecomment},
    stringstyle=\color{codestring},
    keywordstyle=\bfseries,
    aboveskip=6pt,
    belowskip=6pt
}
\lstset{style=shell}

% --- "which machine" banner, used at the top of each chapter -------------
\newcommand{\runson}[1]{%
  \begin{center}
  \fcolorbox{black}{codebg}{\parbox{0.94\textwidth}{\textbf{Runs on:} #1}}
  \end{center}
}

% --- boxed warning, for facts that will cost the reader time if missed ----
\newsavebox{\warnbox}
\newenvironment{warningbox}%
  {\begin{lrbox}{\warnbox}\begin{minipage}{0.90\textwidth}}%
  {\end{minipage}\end{lrbox}%
   \begin{center}\fcolorbox{black}{codebg}{\usebox{\warnbox}}\end{center}}


\title{Carolus / RoboMaster S1 --- Complete Technical Manual}
\author{}
\date{\today}

\begin{document}

\maketitle

\begin{abstract}
This manual is self-sufficient: following it from Chapter~\ref{ch:poweron}
onward requires no other document, inherited tutorial, or vendor PDF. It
covers the complete path from unboxing a rooted RoboMaster S1 and a
Raspberry Pi 4B to a running Carolus/uVGS-2 pose-estimation pipeline with
autonomous beacon docking, using only the corrected, currently-working
configuration of this project --- not the historical or inherited versions
of any script, launch file, or procedure. Every program listed here is the
version that actually runs today.
\end{abstract}

\tableofcontents

\chapter{Introduction}

\section{What this manual covers}

This is the complete, self-contained setup and operation manual for the
Carolus/RoboMaster S1 platform: a DJI RoboMaster S1 ground robot, rooted to
expose the RoboMaster EP SDK, carrying a Raspberry Pi~4B at the turret
position, running Carolus (also called uVGS-2) --- a 4-LED beacon
pose-estimation system --- to autonomously search for, align to, approach,
and dock at a fixed beacon.

Everything in this manual has been executed and verified on this exact
hardware. Where a step has not been physically re-verified on this specific
unit but is inherited from a documented, working precedent (e.g.\ a rooting
procedure that predates this project), that is stated explicitly at the
point it occurs.

\section{Hardware}

\begin{itemize}
  \item \textbf{Robot:} DJI RoboMaster S1, rooted via the
  \texttt{s1\_sdk\_hack\_v0.0.5} community hack, which exposes the
  RoboMaster EP SDK (chassis drive, gimbal control, onboard camera, wheel
  odometry, IMU) over a local API not officially released for the S1.
  \item \textbf{Companion computer:} Raspberry Pi~4B, mounted at the turret,
  connected to the robot over USB via RNDIS (a virtual Ethernet link over
  USB). The Pi holds the single connection to the robot's SDK, runs the
  camera/beacon-tracking node, and --- in the recommended configuration ---
  \textbf{runs Carolus itself}, so the full perception pipeline is onboard
  and the robot needs no network link to estimate its pose.
  \item \textbf{Lab PC:} runs the operator GUI, and serves as a development
  machine for editing and rebuilding. Not physically connected to the robot;
  communicates with the Pi over the local network. \textbf{Carolus can also
  be built and run here} --- both options are documented and compared in
  Section~\ref{sec:where-carolus-runs}. In the recommended configuration it
  runs on the Pi, so the robot does not depend on a network link to the PC to
  localise; running it on the PC is the better choice while developing.
  \item \textbf{Camera:} the S1's own internal camera (built into the
  gimbal/head), streamed over the SDK. This project does not use an external
  Intel RealSense D455 --- some inherited reference material describes a
  D455-based pipeline (see Chapter~\ref{ch:calibration}'s calibration
  discussion); that material is noted where it applies to calibration
  \emph{technique}, not to this robot's actual camera hardware.
\end{itemize}

\section{Software stack}

\begin{itemize}
  \item \textbf{Operating systems:} Raspberry Pi runs Ubuntu 20.04.6~LTS
  (aarch64). The lab PC runs Ubuntu 22.04.5~LTS (x86\_64) --- not the
  officially supported target for ROS Noetic, which requires local
  workarounds documented in Chapter~\ref{ch:carolus-node}.
  \item \textbf{Middleware:} ROS Noetic (1.17.4) on both machines.
  \item \textbf{Robot control:} the official \texttt{robomaster} Python SDK,
  requiring Python~3.7 specifically.
  \item \textbf{Vision:} Carolus/uVGS-2, a C++ ROS node built on OpenCV,
  Eigen3, and Ceres Solver, performing 4-LED blob detection and a
  Perspective-4-Point (P4P) pose solve.
\end{itemize}

\section{What Carolus is}

Carolus (internally named uVGS-2, Unified Visual Guidance System version~2)
is a monocular pose-estimation system originally developed for NASA's
Astrobee free-flying robot aboard the International Space Station, derived
from the earlier SVGS system. It detects four illuminated beacon points
(LEDs) in a camera image, identifies which detected blob corresponds to
which known 3D point on the physical beacon, and solves a Perspective-4-Point
optimization (via Ceres Solver, not OpenCV's \texttt{solvePnP} family) to
recover the beacon's full 6-DOF pose relative to the camera.

\textbf{Carolus's own native coordinate frame is left-handed.} This is not
obvious from the code and is stated here because it matters for every
transform downstream: converting a Carolus pose or orientation into ROS's
right-handed convention (REP-103) requires a true change-of-basis on the
rotation, not a component-wise permutation of a quaternion's vector part ---
the latter is only valid for a proper (determinant~$+1$) rotation, and a
left-to-right-handed conversion is improper (determinant~$-1$). The correct
method (quaternion composition, not permutation) is given at the point it is
needed in this manual.

\section{The four operating states}

Once running, the robot's autonomous behavior is a state machine:

\begin{description}
  \item[SEARCH] The gimbal sweeps without moving the chassis, looking for
  the beacon. If not found from the current heading, the chassis rotates to
  a new heading and the gimbal sweeps again (a ``fan search'' pattern).
  \item[ALIGN] Once the beacon is detected, the chassis rotates in place
  until it faces the beacon.
  \item[APPROACH] The gimbal tracks the beacon while the chassis advances
  along the same heading.
  \item[STOP] Triggered once within a target stopping distance of the
  beacon.
\end{description}

A separate, more deliberate \textbf{docking} behavior performs a calibrated,
verified approach-and-stop sequence rather than the lighter continuous
SEARCH/ALIGN/APPROACH loop. It runs through this project's own operator
tooling, not yet covered in this manual (see the note at the end of
Chapter~\ref{ch:carolus-node}).

\section{How to use this manual}

Follow the chapters in order for a from-scratch setup: power the robot on
and confirm it is rooted (Chapter~\ref{ch:poweron}), root it if it is not
(Chapter~\ref{ch:rooting}), set up the Raspberry Pi
(Chapter~\ref{ch:raspberrypi}), bring up the network link
(Chapter~\ref{ch:network}), bring up the camera pipeline
(Chapter~\ref{ch:camera}), calibrate the camera
(Chapter~\ref{ch:calibration}), then build and launch Carolus
(Chapter~\ref{ch:carolus-node}). At that point the core pipeline is running
and can be driven manually over the topics named in
Chapter~\ref{ch:camera}; the operator GUI and autonomous docking behavior
are this project's own tooling on top of it, documented separately.

Every command, parameter value, and code listing in this manual is the
\emph{current, corrected} version --- the one that actually runs on this
project today.

\chapter{Powering On the Robot}
\label{ch:poweron}

\runson{the \textbf{robot itself} (physical handling). The Raspberry Pi and lab PC are not involved until Chapter~\ref{ch:raspberrypi}.}

\section{Before you start}

\begin{itemize}
  \item The S1's smart battery, charged. Charging is required before first
  use even if the battery shows some charge --- this ``wakes up'' the
  battery.
  \item The S1 charger, plugged into mains (100--240\,V, 50/60\,Hz). A full
  charge takes approximately 1\,h30.
  \item The battery only charges between 5\,\textdegree C and
  45\,\textdegree C.
  \item Do not touch the motors, motor mounting plate, or the inside of the
  Mecanum wheels immediately after powering the robot off --- they can be
  hot.
  \item Do not attempt to charge a battery that has dropped below 1\,V per
  cell; this poses a fire risk.
  \item Use the official DJI charger only.
\end{itemize}

\section{Power-on sequence}

The battery compartment is behind the chassis's rear armor panel, so the
sequence is: open, insert, power on, close.

\begin{enumerate}
  \item Press the rear armor release button to open the chassis's rear
  cover.
  \item Insert the smart battery into its compartment, making sure it is
  firmly seated --- a poorly seated battery gives poor contact and can cause
  the robot to lose battery information.
  \item Press and hold the power button (on the battery itself) for more
  than two seconds to power on.
  \item Close the chassis's rear armor.
\end{enumerate}

The Raspberry Pi at the turret is a separate power domain from the robot
chassis; there is no required power-on order between the two, and the
USB/RNDIS cable connecting them is for the SDK link (Chapter~\ref{ch:network}),
not for powering either device.

\section{Confirming the robot is powered on}

Expected light signals at power-on:

\begin{itemize}
  \item \textbf{Gimbal (head) LED:} blinks cyan, counterclockwise.
  \item \textbf{Chassis LED:} solid cyan.
  \item \textbf{Motion controller LED:} blinks blue slowly --- this means
  the motion controller is working normally.
\end{itemize}

\section{Confirming the robot is rooted}

This project assumes a rooted S1 exposing the RoboMaster EP SDK. Without the
root hack (Chapter~\ref{ch:rooting}), SDK commands (chassis, gimbal, camera)
will not respond, even though the robot otherwise powers on and behaves
normally through the stock RoboMaster app.

\textbf{Audible cue at boot:} a correctly rooted S1 emits two chimes instead
of the usual single chime. The hack persists across power cycles --- it does
not need to be reapplied every time the robot is powered on.

\textbf{Concrete test, once the Pi is connected over RNDIS
(Chapter~\ref{ch:network}):} the camera/beacon node
(\texttt{rm\_cam\_beacon.py}, Chapter~\ref{ch:camera}) opens the one and
only SDK connection this project uses:

\begin{lstlisting}[style=code]
self.ep = robot.Robot()
self.ep.initialize(conn_type="rndis")
\end{lstlisting}

\begin{itemize}
  \item \textbf{Success:} the node's log shows the connection succeeding, and
  the robot responds to \texttt{Twist} commands published on
  \texttt{/carolus/cmd\_vel} (Chapter~\ref{ch:camera}).
  \item \textbf{Failure:} the SDK connection fails or times out. First
  confirm the network link itself (\texttt{ping 192.168.42.2} from the Pi,
  Chapter~\ref{ch:network}). If the network link is fine but the SDK
  connection still fails, redo the power-on and listen carefully for the
  double chime --- a single chime means the root hack is not active.
\end{itemize}

\textbf{Do not use any script other than \texttt{rm\_cam\_beacon.py} to test
this.} Two independent scripts each attempting
\texttt{ep.initialize(conn\_type="rndis")} at the same time silently blocks
robot motion after a few commands, with no error surfaced --- this failure
mode is why this project uses exactly one script holding the SDK connection
at any time (see Chapter~\ref{ch:camera}).

\section{Which camera is mounted}

This project uses the S1's own internal camera, not an external Intel
RealSense D455. Physically: a D455 is an external rectangular camera housing
mounted in addition to the robot's own head; if there is no such external
camera, the internal camera is in use. In software, once
\texttt{rm\_cam\_beacon.py} is running (Chapter~\ref{ch:camera}), the
internal camera feed appears on \texttt{/camera/color/image\_raw}
immediately, with no separate D455 driver (\texttt{rs\_camera.launch})
needed.

\chapter{Rooting the S1}
\label{ch:rooting}

\runson{a \textbf{Windows or macOS computer} connected to the robot by micro-USB. Not the Raspberry Pi, not the Ubuntu lab PC --- the RoboMaster app required for this step does not run on Linux.}

If Chapter~\ref{ch:poweron}'s double-chime check already confirmed the robot
is rooted, this chapter is informational only --- do not re-run it.

\section{What the hack does}

DJI's RoboMaster app exposes a limited, sandboxed Python scripting
environment (``the Lab''). The S1's Python subsystem runs with root
privileges internally, but DJI restricts which modules a Lab script can
import via a whitelist attached to the \texttt{\_\_builtins\_\_} object.
Several of the whitelisted modules (needed for legitimate scripting
features) still hold a reference to the real, unrestricted
\texttt{\_\_builtins\_\_} internally. That reference can be recovered and
used to import \texttt{subprocess}, which can then invoke a shell script
already present on the robot's filesystem (\texttt{adb\_en.sh}) to enable
\texttt{adb} and open a root shell over USB.

Once a root shell is available, a patch is pushed to the robot's filesystem
that installs the RoboMaster EP SDK's full command set (the EP is DJI's
higher-end sibling to the S1; underneath, the two share nearly identical
firmware) onto the S1, which DJI has never officially released to S1 owners.

\textbf{Warning, carried over from the hack's own documentation, to be taken
literally:} this has the potential to break the robot. It did not break this
project's unit, but there is no guarantee for any given unit.

\section{Step-by-step procedure}

\subsection{What you need}

\begin{itemize}
  \item A computer running Windows or macOS (the RoboMaster app requires
  one of these to configure the S1; the Raspberry Pi is not involved in this
  chapter).
  \item Android SDK Platform-Tools (\texttt{adb.exe} and friends) --- unzip
  anywhere, remember the path.
  \item A micro-USB cable.
  \item The RoboMaster app, with a DJI account.
\end{itemize}

\subsection{Procedure}

\begin{enumerate}
  \item Connect the computer to the S1's Intelligent Controller Micro-USB
  port with the micro-USB cable. Ensure no other Android device is connected
  over USB during this procedure.

  \item Power on the robot (Chapter~\ref{ch:poweron}) and connect the
  RoboMaster app to it via its own Wi-Fi hotspot (visible on top of the
  gimbal, with a password shown alongside it).

  \item In the app: \textbf{Lab} $\to$ \textbf{DIY Programming} $\to$ select
  \textbf{Python} on the top bar $\to$ \textbf{New} to create a new program.

  \item Paste the following into the program editor exactly. This is the
  current, working root script for this hack version --- do not use an
  older variant that imports \texttt{'subprocess'} as a literal string or
  reads from \texttt{rm\_log} instead of \texttt{rm\_define}; a firmware
  update closed that specific path.

\begin{lstlisting}[style=code,language=Python]
def root_me(module):
    __import__=rm_define.__dict__['__builtins__']['__import__']
    return __import__(module,globals(),locals(),[],0)

builtins=root_me('builtins')
x=root_me('sub'+'process')
proc=x.Popen('/system/bin/adb_en.sh',shell=True,executable='/system/bin/sh',
             stdout=x.PIPE,stderr=x.PIPE)
\end{lstlisting}

  \item Run the program from within the Lab. A correct run shows no
  compilation errors and the console prints \texttt{Execution Complete}. Do
  not close the RoboMaster app.

  \item On the computer, open a terminal in the folder containing the
  Android Platform-Tools and run:

\begin{lstlisting}[style=shell]
adb devices
\end{lstlisting}

  The S1 should appear in the device list.

  \item Confirm root access:

\begin{lstlisting}[style=shell]
adb shell
\end{lstlisting}

  \item With the root shell open, from a second terminal, in the folder
  containing the hack's distribution files (\texttt{dji\_scratch},
  \texttt{dji.json}, \texttt{dji\_hdvt\_uav}, \texttt{dji\_sys},
  \texttt{patch.sh}), push the patch:

\begin{lstlisting}[style=shell]
adb shell rm -rf /data/dji_scratch
adb push dji_scratch /data/.
adb push dji.json /data/.
adb push dji_hdvt_uav /data/.
adb push dji_sys /data/.
adb shell chmod 755 /data/dji_hdvt_uav
adb shell chmod 755 /data/dji_sys
adb push patch.sh /data/.
\end{lstlisting}

  \item Restart the S1 fully (hold the power button until it turns off,
  then hold it again to turn back on). On boot, listen for \textbf{two
  chimes instead of one} --- this confirms the hack is active
  (Chapter~\ref{ch:poweron}'s check).

  \item Close the RoboMaster app and power off the robot; the setup for
  this chapter is complete.
\end{enumerate}

\section{A pitfall to know before rooting}

Starting with hack version 0.0.3, \texttt{adb} is re-enabled by default in
the patch, which \textbf{disables RNDIS support} --- the exact link this
project's Pi$\leftrightarrow$robot connection depends on
(Chapter~\ref{ch:network}). If RNDIS stops working after rooting, comment
out the \texttt{adb} line in \texttt{patch.sh} before pushing the patch, or
re-push a corrected \texttt{patch.sh} and restart the robot again.

\chapter{Setting Up the Raspberry Pi}
\label{ch:raspberrypi}

\runson{the \textbf{Raspberry Pi 4B} (mounted on the robot), with a keyboard/monitor for first boot or over SSH from the lab PC afterwards.}

This chapter takes a blank Raspberry Pi~4B to a working Ubuntu~20.04
installation with Python~3.7, the RoboMaster SDK, ROS Noetic, and SSH
access. If the Pi is already set up, skip to Chapter~\ref{ch:network}.

\section{Flashing Ubuntu Server}

\begin{enumerate}
  \item Download the Ubuntu Server ARM64 image from
  \url{https://cdimage.ubuntu.com/releases/focal/release/}
  (\texttt{ubuntu-20.04.5-preinstalled-server-arm64+raspi.img.xz}).
  \item Decompress it:
\begin{lstlisting}[style=shell]
unxz ubuntu-20.04.5-preinstalled-server-arm64+raspi.img.xz
\end{lstlisting}
  \item Identify the SD card device (do this carefully --- writing to the
  wrong device destroys data on it):
\begin{lstlisting}[style=shell]
lsblk
\end{lstlisting}
  \item Flash it:
\begin{lstlisting}[style=shell]
sudo dd if=ubuntu-20.04.5-preinstalled-server-arm64+raspi.img of=/dev/sdX bs=4M status=progress conv=fsync
\end{lstlisting}
  \item Boot the Pi with the card inserted, plus keyboard, mouse, display,
  and an Ethernet cable for first boot. Default credentials:
  \texttt{ubuntu}/\texttt{ubuntu} (you will be asked to change the password
  on first login).
\end{enumerate}

\section{Desktop environment (optional)}

Only needed if you want a graphical session on the Pi itself (not required
for this project's normal operation, which runs headless):

\begin{lstlisting}[style=shell]
sudo apt update
sudo apt install ubuntu-desktop
sudo reboot
\end{lstlisting}

\section{Enabling headless SSH access}

If SSH was not already enabled at flash time (Raspberry Pi Imager's advanced
options support this directly), enable it on first boot:

\begin{lstlisting}[style=shell]
sudo systemctl enable ssh
sudo systemctl start ssh
sudo systemctl status ssh   # optional: confirm it's running
\end{lstlisting}

Find the Pi's IP address on the local network:

\begin{lstlisting}[style=shell]
hostname -I
\end{lstlisting}

Connect from the client machine, first with a password, then switch to
key-based login:

\begin{lstlisting}[style=shell]
# First connection (password):
ssh <username>@<pi_ip>

# On the client: generate a key
ssh-keygen -t ed25519 -C "client-key"

# Copy it to the Pi
ssh-copy-id <username>@<pi_ip>

# Test key-based login (no password prompt expected)
ssh <username>@<pi_ip>
\end{lstlisting}

A convenient alias, in \texttt{\textasciitilde/.ssh/config} on the client:

\begin{lstlisting}[style=code]
Host raspberry
   HostName <pi_ip>
   User <username>
   IdentityFile ~/.ssh/id_ed25519
\end{lstlisting}

\textbf{Note the exact directive name: \texttt{IdentityFile}, not
\texttt{IdentifyFile}} --- a common typo that OpenSSH silently rejects the
config over.

Connect with \texttt{ssh raspberry} from then on.

\textbf{Optional hardening, once key-based login is confirmed working:}

\begin{lstlisting}[style=shell]
# Firewall
sudo apt install -y ufw
sudo ufw allow OpenSSH
sudo ufw enable

# Disable password authentication (only after keys work!)
sudo nano /etc/ssh/sshd_config
# Ensure: PasswordAuthentication no
#         PubkeyAuthentication yes
sudo systemctl restart ssh
\end{lstlisting}

\section{Building Python 3.7 from source}

The RoboMaster SDK requires Python~3.7 specifically; Ubuntu~20.04 ships a
newer default Python, so 3.7 is built from source into a separate location
rather than replacing the system Python.

\begin{lstlisting}[style=shell]
sudo apt update
sudo apt install -y build-essential libssl-dev zlib1g-dev \
    libreadline-dev libsqlite3-dev libbz2-dev libexpat1-dev \
    liblzma-dev tk-dev libffi-dev uuid-dev

cd ~/Python-3.7.17/
./configure --enable-optimizations
make -j4
sudo make altinstall
\end{lstlisting}

Using \texttt{altinstall} (not \texttt{install}) is deliberate: it installs
as \texttt{python3.7} alongside the system Python, without overwriting it.

Confirm the SSL module built correctly (a common source of a broken venv
later if skipped):

\begin{lstlisting}[style=shell]
python3.7 -c "import ssl; print(ssl.OPENSSL_VERSION)"
\end{lstlisting}

\section{Virtual environment and the RoboMaster SDK}

\begin{lstlisting}[style=shell]
python3.7 -m venv env
source env/bin/activate
pip install --upgrade pip
pip install robomaster numpy
pip install "opencv-python==4.5.5.64"
\end{lstlisting}

The pinned OpenCV version is a historically validated compatible version for
this SDK/Python combination on this hardware; it is not a hard requirement
enforced elsewhere in this manual, and later versions have since been used
successfully on the lab PC (Chapter~\ref{ch:carolus-node}) without issue.

\section{Installing ROS Noetic}

\begin{lstlisting}[style=shell]
sudo sh -c 'echo "deb http://packages.ros.org/ros/ubuntu \
$(lsb_release -sc) main" > /etc/apt/sources.list.d/ros-latest.list'
sudo apt install curl
curl -s https://raw.githubusercontent.com/ros/rosdistro/master/ros.asc \
    | sudo apt-key add -
sudo apt update

# Full desktop version:
sudo apt install ros-noetic-desktop-full
# --- or, if the Pi struggles with the desktop version's resource use ---
sudo apt install ros-noetic-ros-base

echo "source /opt/ros/noetic/setup.bash" >> ~/.bashrc
source ~/.bashrc

sudo apt install python3-rosdep python3-rosinstall \
    python3-rosinstall-generator python3-wstool build-essential
sudo rosdep init
rosdep update
\end{lstlisting}

On this project's Pi, the lighter \texttt{ros-noetic-ros-base} is what
actually runs day to day --- the Pi only needs to run
\texttt{rm\_cam\_beacon.py} (Chapter~\ref{ch:camera}), not a full desktop
ROS toolchain.

\section{Getting this project's code onto the Pi}
\label{sec:pi-code}

The Pi needs this project's own workspace --- specifically
\texttt{carolus\_ws/src/robomaster\_cam/}, which holds
\texttt{rm\_cam\_beacon.py} (Chapter~\ref{ch:camera}) and its companion
scripts. Clone the project's repository directly onto the Pi:

\begin{lstlisting}[style=shell]
git clone git@github.com:Arracheur2Bonnet/robotmaster.git carolus_ws
\end{lstlisting}

This uses an SSH key registered against \emph{your own} GitHub account
(Chapter~\ref{ch:carolus-node} covers generating and registering one). If
the repository is private you will also need collaborator access --- ask the
project supervisor. Either way, \textbf{do not paste a shared personal
access token into the clone URL}: it grants whoever reads it the full rights
of the account that made it, and cannot be traced back to an individual. If
SSH is not set up, an HTTPS clone with your own GitHub credentials works the
same way:

\begin{lstlisting}[style=shell]
git clone https://github.com/Arracheur2Bonnet/robotmaster.git carolus_ws
\end{lstlisting}

No build step is needed on the Pi side for \texttt{robomaster\_cam} ---
it is a pure-Python package, run directly (Chapter~\ref{ch:camera}).

\section{A pitfall to expect at first launch}

The RNDIS network interface (Chapter~\ref{ch:network}) does not come up
automatically at Pi boot, even once everything above is configured
correctly. This is expected --- Chapter~\ref{ch:network} gives the exact
two commands to run at the start of every session.

\chapter{Network Setup (RNDIS)}
\label{ch:network}

\runson{the \textbf{Raspberry Pi} for the interface commands; one step at the end runs on the \textbf{lab PC} and is marked as such.}

The robot exposes an RNDIS virtual Ethernet interface over the same USB
cable used for the SDK connection. The robot holds a fixed address,
\texttt{192.168.42.2}; the Raspberry Pi must be configured on the same
subnet.

\section{Bringing up the link}

\textbf{Run these two commands at the start of every session, on the Pi.}
The interface does not persist across Pi reboots --- this is expected
behavior on this hardware, not a fault, and there is no step elsewhere in
this manual that avoids needing it.

\begin{lstlisting}[style=shell]
sudo ip link set eth1 up
sudo ip addr add 192.168.42.3/24 dev eth1
\end{lstlisting}

If the interface is not named \texttt{eth1} on your system, identify it
first:

\begin{lstlisting}[style=shell]
ip -br addr
# look for eth1 / usb0 / enx<MAC> -- whichever one is meant to carry 192.168.42.x
\end{lstlisting}

\textbf{Never assign \texttt{192.168.42.2} to the Pi} --- that address
belongs to the robot.

\subsection{If you intend to automate this over SSH}

Typing the two commands by hand prompts for a password, which is fine
interactively. \textbf{Any script that brings the interface up over SSH needs
passwordless \texttt{sudo} on the Pi}, otherwise it hangs silently waiting
for a prompt it cannot display --- the process simply appears frozen, with no
error. Check before automating anything:

\begin{lstlisting}[style=shell]
ssh <user>@<pi> 'sudo -n true && echo "passwordless sudo OK"'
\end{lstlisting}

If it fails, grant it with \texttt{visudo} on the Pi (scope it to the two
\texttt{ip} commands rather than granting blanket rights).

\section{Verifying the link}

\begin{lstlisting}[style=shell]
ping -c 4 192.168.42.2
\end{lstlisting}

Expect 0\% packet loss, sub-millisecond round-trip time once the interface
is correctly up.

\section{A quick automated alternative}

Rather than typing the two commands above by hand every session, they can be
scripted:

\begin{lstlisting}[style=shell]
sudo ip link set eth1 up && sudo dhclient eth1
\end{lstlisting}

This project's own operator tooling (not covered in this manual) runs the
equivalent of the first two-command form automatically as part of its first
launch step, over SSH from the lab PC.

\section{Lab PC hostname resolution}

If any ROS tool on the lab PC (e.g.\ \texttt{rqt\_image\_view}) shows a
black image or errors with something like
\texttt{Couldn't find AF\_INET address for [ubuntu]}, add the Pi's hostname
to \texttt{/etc/hosts} on the lab PC:

\begin{lstlisting}[style=shell]
echo "192.168.0.103 ubuntu" | sudo tee -a /etc/hosts
\end{lstlisting}

Unlike the RNDIS interface on the Pi, this fix is permanent on the lab PC
and does not need to be repeated.

\section{Summary of addresses used in this manual}

\begin{center}
\begin{tabular}{ll}
\toprule
Device & Address \\
\midrule
Robot (over RNDIS) & \texttt{192.168.42.2} \\
Raspberry Pi (over RNDIS, to the robot) & \texttt{192.168.42.3} \\
Raspberry Pi (on the local network, to the lab PC) & \texttt{192.168.0.103} \\
\bottomrule
\end{tabular}
\end{center}

The last of these is a DHCP-assigned local-network address and may differ on
a different network; find it with \texttt{hostname -I} run on the Pi
(Chapter~\ref{ch:raspberrypi}).

\chapter{The Camera and Robot Bridge}
\label{ch:camera}

\runson{the \textbf{Raspberry Pi} --- this is the only machine that ever connects to the robot SDK. The verification commands at the end run from the \textbf{lab PC}.}

\section{Design: a single SDK connection}

The Pi runs exactly one script that ever connects to the robot's SDK:
\texttt{rm\_cam\_beacon.py} (in
\texttt{carolus\_ws/src/robomaster\_cam/scripts/}). It is not just a camera
capture script --- it is this project's single bridge to the entire SDK:
camera stream, chassis drive, gimbal control, odometry, and the IMU
subscription all go through this one script and its one connection.

\textbf{This is a deliberate design choice, not an incidental one.} Two
independent scripts each holding their own
\texttt{ep.initialize(conn\_type="rndis")} connection at the same time
silently blocks robot motion commands after a few calls, with no error
surfaced. \textbf{Never run any other script that calls
\texttt{robot.Robot().initialize(...)} while \texttt{rm\_cam\_beacon.py} is
running.}

\section{Starting the bridge}

On the Pi, with the RNDIS link up (Chapter~\ref{ch:network}) and the Python
3.7 virtual environment active (Chapter~\ref{ch:raspberrypi}):

\begin{lstlisting}[style=shell]
source ~/env/bin/activate
python3.7 carolus_ws/src/robomaster_cam/scripts/rm_cam_beacon.py
\end{lstlisting}

In normal operation this is launched automatically by this project's own
operator tooling, over SSH from the lab PC --- it is not usually started by
hand.

\section{What it connects to}

\begin{lstlisting}[style=code,language=Python]
self.ep = robot.Robot()
self.ep.initialize(conn_type="rndis")
self.chassis = self.ep.chassis
self.gimbal  = self.ep.gimbal
self.cam     = self.ep.camera
self.cam.start_video_stream(display=False)
\end{lstlisting}

From there it subscribes to the SDK's own push-data channels
(\texttt{sub\_position}, \texttt{sub\_attitude}, \texttt{sub\_velocity},
\texttt{sub\_imu}, \texttt{sub\_angle} for the gimbal), converts each
callback into a ROS message, and publishes:

\begin{center}
\begin{tabular}{ll}
\toprule
Topic & Content \\
\midrule
\texttt{/camera/color/image\_raw} & Camera frames (\texttt{sensor\_msgs/Image}) \\
\texttt{/odom} & Wheel odometry, republished from the SDK \\
\texttt{/imu} & Raw IMU, 50\,Hz requested --- \textbf{units: see warning} \\
\texttt{/carolus/gimbal\_yaw\_rel} & Gimbal yaw relative to the chassis \\
\texttt{/carolus/gimbal\_yaw\_ground} & Gimbal yaw in the world frame (absolute) \\
\bottomrule
\end{tabular}
\end{center}

\begin{warningbox}
\textbf{\texttt{/imu} units are \emph{g}, not m/s\textsuperscript{2}
(measured 2026-08-04).} At rest on a flat floor, this topic reports
\texttt{linear\_acceleration.z = -1.00664}. The SDK returns accelerations in
\emph{g}, and \texttt{rm\_cam\_beacon.py} republishes them unconverted ---
whereas \texttt{sensor\_msgs/Imu} mandates
m/s\textsuperscript{2}. Any consumer of this topic therefore sees
accelerations $9.81\times$ too small under a message contract that says
otherwise, and neither Kalibr nor MINS will detect it. The sign is also
suspect: a REP-103 $z$-up frame at rest should read $+9.81$, not $-1.0$\,g.
Gyroscope units and axis signs are likewise unconfirmed --- a reading at rest
cannot distinguish rad/s from \textdegree/s when the true rate is zero.

\textbf{Before using this topic}, verify on your own robot: gravity should
read $\approx\!9.81$ on exactly one axis at rest, and a known slow rotation
should give the expected sign and magnitude on the corresponding gyro axis.
Fix scale, sign and gyro units together --- correcting the scale alone yields
a topic that looks right while its frame convention is still unverified.

Two further constraints: this SDK accepts only 1, 5, 10, 20 and 50\,Hz for
this subscription, so 50\,Hz is the ceiling --- below the 200--500\,Hz that
camera--IMU calibration tools such as Kalibr
(Chapter~\ref{ch:calibration}) normally expect. The frequency is exposed as
the \texttt{RM\_IMU\_FREQ} environment variable so the other legal values can
be tried without editing the script.
\end{warningbox}

\begin{warningbox}
\textbf{If \texttt{/imu} is silent, check the publisher creation order
before suspecting the robot.} Every ROS publisher written to from an SDK
callback must be created \emph{before} the corresponding
\texttt{sub\_*} subscription. The robot begins pushing data the moment the
subscription is registered, and on the SDK build used here the callback runs
directly in the DDS dispatcher thread: an \texttt{AttributeError} on a
not-yet-created publisher kills that thread, after which \emph{no} telemetry
callback ever runs again. The symptom is indistinguishable from a robot that
sends nothing --- the topic is advertised, \texttt{sub\_imu} returns
\texttt{True}, and ROS logs nothing. This cost a full session to diagnose.
\end{warningbox}

and subscribes to a small set of control topics used by this project's own
higher-level tooling (not covered in this manual):
\texttt{/pose} (Carolus's own output, consumed here to drive autonomous
behavior), \texttt{/carolus/mode}, \texttt{/carolus/cmd\_vel},
\texttt{/carolus/gimbal\_vel}, \texttt{/carolus/wheels},
\texttt{/carolus/gimbal\_lock}, \texttt{/carolus/gimbal\_lock\_period}, and
\texttt{/carolus/gimbal\_recenter}.

\section{Verifying it is running}

From the lab PC, with the ROS master reachable:

\begin{lstlisting}[style=shell]
rostopic list | grep camera
rostopic hz /camera/color/image_raw
rostopic echo -n1 /odom
rostopic echo -n1 /imu
\end{lstlisting}

Measured on the Raspberry Pi with the robot connected (2026-08-04):
\texttt{/camera/color/image\_raw} at \textbf{30.1\,Hz},
\texttt{/odom} at \textbf{15.4\,Hz}, \texttt{/imu} at \textbf{47.3\,Hz}
against a 50\,Hz subscription. Expect at least one non-empty message on
\texttt{/odom} and \texttt{/imu}; if \texttt{/imu} hangs, read the second
warning above before concluding anything about the robot.
If the camera topic is silent, confirm the RNDIS link
(Chapter~\ref{ch:network}) before suspecting the script itself --- a dead
network link and a dead camera stream look identical from this check alone.

\section{Gimbal yaw: two signals, two different meanings}

\texttt{/carolus/gimbal\_yaw\_rel} and \texttt{/carolus/gimbal\_yaw\_ground}
are not interchangeable, and confusing them is a documented source of bugs
in this project's history:

\begin{itemize}
  \item \textbf{\texttt{gimbal\_yaw\_ground}} is the gimbal's heading in the
  world frame. The S1's gimbal is yaw-stabilized: when the chassis rotates
  without an explicit gimbal command, the gimbal holds its absolute heading.
  This value stays essentially constant across a chassis rotation.
  \item \textbf{\texttt{gimbal\_yaw\_rel}} is the gimbal's angle
  \emph{relative to the chassis}. Because the gimbal holds its world heading
  while the chassis turns underneath it, this value \emph{absorbs} the
  chassis's rotation --- it is not a stable ``how far off is the beacon''
  signal across a chassis maneuver, it is a signal that actively tracks how
  much the chassis has turned since the gimbal last held still relative to
  it.
\end{itemize}

Use \texttt{gimbal\_yaw\_ground} when you need to know the camera's absolute
heading in the world. Use \texttt{gimbal\_yaw\_rel} only for what it
actually is: the angle between the gimbal and the chassis at this instant,
not a proxy for beacon bearing across a chassis rotation.

\section{The internal camera vs.\ an external D455}

This robot uses its own internal camera exclusively. If, when following any
inherited reference material, you encounter instructions to launch
\texttt{rs\_camera.launch} or install \texttt{librealsense}, that material
concerns a different camera than the one on this robot and does not apply
here --- see the calibration discussion in Chapter~\ref{ch:calibration} for
where that material remains genuinely useful (the calibration
\emph{technique}, not the camera driver).

\chapter{Camera Calibration}
\label{ch:calibration}

\runson{the \textbf{lab PC} (or any machine with MATLAB / a Kalibr install). Images and rosbags come from the robot camera over the network.}

\section{Current values in use}

The intrinsics and distortion coefficients currently loaded by
\texttt{testcarolus.launch} (Chapter~\ref{ch:carolus-node}) are:

\begin{center}
\begin{tabular}{ll}
\toprule
Parameter & Value \\
\midrule
$f_x$ & 546.1957 \\
$f_y$ & 547.0838 \\
$c_x$ & 575.6041 \\
$c_y$ & 372.1876 \\
Distortion (plumb\_bob: $k_1,k_2,p_1,p_2,k_3$) & $-0.1479,\ 0.1452,\ -0.0023,\ 0.0025,\ 0.0070$ \\
\bottomrule
\end{tabular}
\end{center}

These come from \texttt{robomaster\_cam/config/camera\_info.yaml}
(\texttt{camera\_name: robomaster\_s1}, $1280\times720$), validated
pragmatically --- Carolus produces a stable, coherent pose with these values
in practice. They are \textbf{not} the result of a formal metric
calibration campaign (checkerboard/AprilGrid target, dedicated calibration
tool run to convergence with a reported reprojection error). If you are
setting up a new unit, or want a rigorous calibration, follow the procedure
below rather than reusing these values --- they are specific to this
project's camera and are not guaranteed to transfer to a different physical
unit.

\section{Recommended method: a checkerboard with the MATLAB Camera
Calibration Toolbox}

This procedure has not yet been run against this project's own camera ---
the values in the previous section are what is actually loaded today. It is
included here as the recommended path for whoever runs the next
calibration, based on a direct comparison against Kalibr on a comparable
camera (Section~\ref{sec:kalibr-alt}).

\begin{enumerate}
  \item Print a checkerboard target with known square size (25\,mm squares
  is a proven working size for this class of camera).
  \item Capture a set of images of the board with the robot's camera
  (Chapter~\ref{ch:camera}): at least 10 images, 15 or more for better
  accuracy. The board should occupy roughly 20--25\% of the frame area, and
  be tilted more than 45\textdegree{} from the camera's optical axis in each
  shot to avoid detection errors. Capture from a wide variety of positions
  and orientations across the field of view.
  \item In MATLAB (desktop or MATLAB Online, no additional software
  required): open the \textbf{Camera Calibration app} (Apps tab), add the
  captured images, select the checkerboard pattern with the correct square
  size, and enable \textbf{tangential distortion} and \textbf{three radial
  distortion coefficients} (this camera class shows moderate wide-angle
  distortion that the default two-coefficient model under-fits).
  \item Run the calibration. Export the resulting coefficients to the MATLAB
  workspace and generate the calibration script --- this script contains the
  intrinsic matrix, distortion coefficients, and reprojection errors in a
  form that can be transcribed directly into \texttt{testcarolus.launch}'s
  \texttt{fx}/\texttt{fy}/\texttt{cx}/\texttt{cy}/\texttt{distortion}
  parameters (Chapter~\ref{ch:carolus-node}).
\end{enumerate}

\section{Alternative / complementary method: Kalibr}
\label{sec:kalibr-alt}

Kalibr (ETH Zürich) is the standard tool if a camera--IMU joint calibration
is needed (required for any future MINS/visual-inertial integration, not
needed for Carolus's own pose estimation, which does not use the IMU).

\textbf{Not executed here.} The camera--IMU half of this procedure has never
been run on this robot. The camera-only calibration below is unaffected and
is the part that matters for Carolus. Two constraints to know before planning
a camera--IMU session on this platform: the SDK caps the IMU subscription at
50\,Hz, well below the 200--500\,Hz Kalibr expects for a well-conditioned
solution; and \texttt{/imu} publishes accelerations in \emph{g} rather than
the m/s\textsuperscript{2} the message type mandates
(Chapter~\ref{ch:camera}), with gyroscope units and axis signs still
unconfirmed. Correct the units and verify the frame first --- Kalibr will not
detect either error, and every noise parameter it produces would be wrong by
the same silent factor.

\subsection{Target}

A $6\times6$ AprilGrid, tag size $0.089\,\mathrm{m}$, tag spacing $0.3\times$
tag size (an $\sim\!0.8\times0.8\,\mathrm{m}$ printed board). Measure the
actual printed square size and use the measured value, not the nominal one
--- even a $\sim\!1\%$ sizing error biases the recovered scale.

\begin{lstlisting}[style=code]
target_type: 'aprilgrid'
tagCols: 6
tagRows: 6
tagSize: 0.089        # metres -- MEASURE your own printed board
tagSpacing: 0.3        # ratio of gap to tag size
codeOffset: 0
\end{lstlisting}

\subsection{Installing Kalibr}

\begin{lstlisting}[style=shell]
cd ~/camera_ws/src
git clone https://github.com/ethz-asl/kalibr.git
cd ~/camera_ws
rosdep install --from-paths src --ignore-src -y
sudo apt install python3-rosdep python3-rosinstall python3-rosinstall-generator \
    build-essential python3-catkin-tools
sudo apt-get install libv4l-dev doxygen python3-igraph
sudo apt install python3-pip
pip install wxPython   # can take a long time -- be patient
catkin_make -DCMAKE_BUILD_TYPE=Release
source ~/camera_ws/devel/setup.bash
\end{lstlisting}

\textbf{Do not separately clone \texttt{numpy\_eigen} or
\texttt{Schweizer-Messer}} --- Kalibr bundles its own copies, and cloning
them independently into \texttt{src/} produces duplicate-package build
errors. If you hit \texttt{ModuleNotFoundError: No module named 'sm'} after
building, verify with:

\begin{lstlisting}[style=shell]
python3 -c "import sm; print(sm.__file__)"
\end{lstlisting}

and if it fails, remove any independently-cloned copy of
\texttt{Schweizer-Messer}, then rebuild:

\begin{lstlisting}[style=shell]
rm -rf build/ devel/ install/
catkin_make -DCMAKE_BUILD_TYPE=Release
\end{lstlisting}

\subsection{Capturing data and running the calibration}

Record a rosbag while moving the camera (or the robot itself, for a joint
camera--IMU calibration) through the AprilGrid's field of view for 60--90
seconds, exciting all six degrees of freedom (translate near/far/side to
side, rotate about all three axes, not just yaw).

\begin{lstlisting}[style=shell]
rosbag record -O calib_data.bag /camera/color/image_raw
\end{lstlisting}

Run the calibration:

\begin{lstlisting}[style=shell]
export KALIBR_MANUAL_FOCAL_LENGTH_INIT=1
kalibr_calibrate_cameras \
    --bag calib_data.bag --target target.yaml \
    --models pinhole-radtan \
    --topics /camera/color/image_raw \
    --approx-sync 0.01
\end{lstlisting}

If Kalibr reports \texttt{Initialization of focal length failed}, provide a
manual seed value (\texttt{600} is a reasonable starting guess for a
$640\times480$-class image; scale proportionally for other resolutions).

\subsection{Reading the result}

Kalibr writes three files: a \texttt{.txt} summary, a \texttt{.pdf} of
reprojection-error plots, and a \texttt{.yaml} with the calibration
parameters. \textbf{Target: reprojection error under 0.2--0.5 pixels.} In
practice, a completed run can land a little outside that target (observed:
0.59--0.67\,px std on a comparable camera) without necessarily being
unusable --- treat the target as a goal to aim for, not a hard pass/fail
gate, and judge the calibration by whether downstream pose estimates are
stable and physically plausible.

\section{Camera--IMU time offset}

If a Kalibr camera--IMU joint calibration is run, its output
\texttt{timeshift\_cam\_imu} is in \textbf{seconds}, with the convention
$t_{\mathrm{imu}} = t_{\mathrm{cam}} + \mathrm{shift}$. This matches
OpenVINS's own internal convention exactly --- the value can be used
directly with no sign flip if it is later consumed by an OpenVINS-based
filter.

\chapter{Building and Launching Carolus}
\label{ch:carolus-node}

\runson{\textbf{either machine} --- both builds are documented here. The
\textbf{Raspberry Pi 4B} is the recommended target for operation, the
\textbf{lab PC} the better one for development;
Section~\ref{sec:where-carolus-runs} compares them with the measurements
behind the recommendation.}

Carolus (the ROS package name is \texttt{carolus\_astrobee}) runs on
whichever machine you build it on. It has no SDK connection of its own ---
it consumes \texttt{/camera/color/image\_raw} (published by
\texttt{rm\_cam\_beacon.py}, Chapter~\ref{ch:camera}) and publishes a pose on
\texttt{/pose} --- so it only needs to reach the ROS master, not the robot.
That is exactly why it can sit on either side, and why the choice is an
architectural one rather than a technical constraint.

\section{Where Carolus runs}
\label{sec:where-carolus-runs}

Both machines work, and this manual documents both. \textbf{The
recommendation is the Pi}, but the reasons matter more than the verdict, so
they are given with their evidence and their limits rather than as a slogan.

\subsection{The two options, side by side}

\begin{center}
\begin{tabular}{p{0.30\textwidth}p{0.30\textwidth}p{0.30\textwidth}}
\toprule
 & \textbf{Raspberry Pi 4B (recommended)} & \textbf{Lab PC} \\
\midrule
Robot autonomy
 & Complete. No network link is needed to estimate pose.
 & None. Every pose depends on a live link; the robot cannot leave the PC's
   network reach. \\
\addlinespace
Camera input rate (measured)
 & \textbf{30.1\,Hz} --- frames never leave the machine that captured them.
 & 10--15\,Hz --- every frame is an \emph{uncompressed}
   \texttt{sensor\_msgs/Image} at $1280\times720$ ($\approx$2.76\,MB)
   crossing the network first. \\
\addlinespace
Build complexity
 & Ubuntu 20.04 = ROS Noetic's officially supported target. Standard
   \texttt{catkin\_make}, no workarounds (Section~\ref{sec:build2004}).
 & On Ubuntu 22.04: four build-time fixes plus one runtime fix
   (Section~\ref{sec:build2204}). On 20.04, identical to the Pi. \\
\addlinespace
Development comfort
 & Editing and rebuilding C++ over SSH is slower.
 & \textbf{Better.} A workstation compiles faster and edits more
   comfortably; GUI tooling runs here. \\
\addlinespace
Visual-inertial work (MINS/OpenVINS)
 & \textbf{Inside} the upstream tested set (Ubuntu 18/20, Melodic/Noetic).
 & On 22.04, \textbf{outside} it. \\
\bottomrule
\end{tabular}
\end{center}

\subsection{Measured, 2026-08-04 --- and what the numbers do \emph{not} say}
\label{sec:pi-benchmark}

Same robot, same beacon at 1.00\,m, same afternoon --- Carolus run once on
each machine and \texttt{/pose} measured both times:

\begin{center}
\begin{tabular}{lll}
\toprule
 & Carolus on the \textbf{lab PC} & Carolus on the \textbf{Pi} \\
\midrule
\textbf{\texttt{/pose}}, beacon in view & 2.19\,Hz & \textbf{13.04\,Hz} \\
\texttt{/postprocessed/image} & --- & 14.35\,Hz \\
\texttt{/camera} at the Pi's publisher & 30.05\,Hz & 16.52\,Hz \\
Pi CPU & 52--61\% idle & \textbf{12.3\% idle}, load 5.37/4 \\
\bottomrule
\end{tabular}
\end{center}

\textbf{A factor of 5.95.} On the PC every frame is an \emph{uncompressed}
$1280\times720$ \texttt{sensor\_msgs/Image} ($\approx$2.76\,MB) crossing the
network before Carolus can touch it; on the Pi the frames never leave the
machine that captured them.

The other rows say where the limit moved to. \texttt{/pose} 13.0\,Hz,
\texttt{/postprocessed} 14.4\,Hz and the camera at 16.5\,Hz are nearly the
same number --- Carolus consumes almost every frame it is handed. And the
camera drops from 30.05\,Hz to 16.52\,Hz as soon as Carolus runs beside it,
with the Pi at 12.3\% idle: \textbf{on the Pi the bottleneck is now CPU, and
the camera bridge is competing with Carolus for it.} If more pose rate is
needed, that is where to look --- publishing compressed frames, or dropping
the capture resolution.

\begin{warningbox}
\textbf{Verify which machine you are measuring.} Both numbers above are
\texttt{/pose} rates with the same beacon, and the only difference is where
the node ran. It is easy to get this wrong: the ROS master lives on the Pi in
both cases, so a Carolus running on the workstation still looks Pi-side from
the topic list. Confirm with \texttt{pgrep -f carolus\_astrobee} \emph{on the
Pi} before attributing a rate to it.

Two older figures circulate and neither is a pose rate: a $\sim$24\,frames/s
figure for the Pi was counted \textbf{with no beacon in view}, so no P4P
solve ran on any frame, and a $\sim$2.5\,frames/s figure for the lab PC used
a different counting method again. Use the table above, not those.
\end{warningbox}

The Pi already had Ceres 1.14, Eigen 3.3.7 and OpenCV 4.2 installed;
\texttt{catkin\_make} compiled and linked directly on \texttt{aarch64}.
Available memory (7.6\,GB) was never a constraint.

\subsection{So which one, and why}

\textbf{Choose the Pi for anything resembling operation.} The decisive
argument is autonomy, not speed: a robot whose pose estimator lives on
another machine stops working the moment the link does, and that is a
property of the architecture rather than a number that a future optimisation
might change. The 2--3$\times$ camera input rate and the absence of build
workarounds are real secondary benefits.

\textbf{Choose the lab PC for development.} Compiling and iterating on C++ is
faster on a workstation than over SSH, and the operator GUI runs there
anyway. If the PC runs Ubuntu 20.04 the build is identical to the Pi's; on
22.04 it costs the five fixes in Section~\ref{sec:build2204}.

\textbf{A note on visual-inertial navigation.} MINS has been run on this
Raspberry Pi, both on its own simulation and on data recorded from this
robot. It builds on \texttt{aarch64} and produces a \emph{consistent}
estimate on the simulation (RMSE 0.113\textdegree{} / 0.082\,m, NEES 0.9 and
1.4), with no divergence, at \textbf{0.3--0.4$\times$ real time} --- under a
sensor load far heavier than this robot's (IMU at 200\,Hz, \emph{two}
cameras, and a LIDAR).

On a 60\,s recording of \emph{this} robot's \texttt{/imu} and
\texttt{/camera/color/image\_raw}, it processed the whole bag at
\textbf{0.7$\times$ real time}: CAM 118\,ms, IMU 3\,ms, 122\,ms of compute
per 88\,ms of data. So \textbf{``MINS does not run on a Raspberry Pi'' is not
correct} --- it runs, on this robot's own data. What is correct is that
\textbf{it does not reach real time here}, even with one camera, no LIDAR and
a 50\,Hz IMU, and that essentially all of the cost is the camera update. The
0.7$\times$ is also optimistic: \texttt{rosbag record} throttled the sources
to 16\,Hz of camera against 30\,Hz live.

\textbf{Two practical points if you repeat this.} MINS's live-subscribe entry
point is \textbf{ROS 2 only} --- in \texttt{run\_subscribe.cpp} the
\texttt{ROS\_AVAILABLE == 1} branch contains two \texttt{\#include} lines and
nothing else, so under ROS Noetic the binary initialises, subscribes to
nothing, and exits immediately even with data flowing. Use the \texttt{bag}
binary and a recorded rosbag, which is the path \texttt{subscribe.launch}'s
own hard-wired \texttt{rosbag play} node implies. And do not quote pose
accuracy from such a run unless the camera--IMU extrinsic has been
calibrated and the IMU units corrected (Chapter~\ref{ch:camera}) --- with a
placeholder extrinsic and accelerations in \emph{g}, the filter diverges
grossly while the timing figures stay valid.
\section{Getting this project's code onto the machine that will build it}

Clone the project's repository onto whichever machine will build Carolus
--- \textbf{normally the Raspberry Pi} (Section~\ref{sec:where-carolus-runs}),
over SSH from the lab PC (Chapter~\ref{ch:raspberrypi}). Cloning into a
directory named \texttt{carolus\_ws} gives a standard catkin workspace
layout, with the ROS packages under \texttt{src/}:
\texttt{src/carolus\_node/} (the launch files and configuration),
\texttt{src/libuvgs\_astrobee/} (Carolus's own C++ source, already patched
for this robot), \texttt{src/ff\_msgs/} (a message-definition dependency),
\texttt{src/robomaster\_cam/} (the Pi-side bridge, Chapter~\ref{ch:camera}),
plus \texttt{cmake\_shims/} (Ubuntu 22.04 build fixes, see below) and
\texttt{shortcuts/} (this project's operator tooling, not covered in this
manual) at the top level:

\begin{lstlisting}[style=shell]
git clone git@github.com:Arracheur2Bonnet/robotmaster.git carolus_ws
# or, over HTTPS with your own GitHub credentials:
# git clone https://github.com/Arracheur2Bonnet/robotmaster.git carolus_ws
\end{lstlisting}

Repository web address:
\url{https://github.com/Arracheur2Bonnet/robotmaster}

\subsection{If the clone fails, or the link shows 404}

The repository may be private. \textbf{If you open the link above and get
404 Not Found}, that does not mean the address is wrong --- GitHub answers
404 rather than ``403 Forbidden'' for a private repository you have no
access to, deliberately, so that the existence of private repositories is
not confirmed to strangers. A 404 here reads as ``you do not have access
yet''.

To get access, ask the repository owner or the project supervisor to add
your GitHub account as a collaborator (\emph{Settings $\to$ Collaborators
$\to$ Add people}). You will need to give them your GitHub username, and
you must accept the emailed invitation before cloning works.

If the repository is public, both clone commands above work as-is with no
access request --- but the SSH setup below is still worth doing, since it
avoids retyping credentials on every push and is required if you are ever
given write access.

\subsection{Setting up an SSH key}

Cloning over SSH (the first command above) needs an SSH key registered
against \emph{your own} GitHub account. This is a one-time setup per
machine, and it is what replaces the practice --- found in some older
inherited documents --- of pasting a shared personal access token directly
into a clone command. A shared token grants whoever holds it the full
access rights of the account that created it, cannot be attributed to an
individual, and stays valid in every copy of the document that carries it.

\begin{enumerate}
  \item Generate a key on the machine that will clone the repository (the
  Raspberry Pi and the lab PC each need their own):

\begin{lstlisting}[style=shell]
ssh-keygen -t ed25519 -C "your_email@example.com"
# Press Enter to accept the default path (~/.ssh/id_ed25519).
# A passphrase is optional but recommended.
\end{lstlisting}

  \item Display the \textbf{public} half of the key --- this is the only
  part that ever leaves the machine. Never share the file without the
  \texttt{.pub} extension:

\begin{lstlisting}[style=shell]
cat ~/.ssh/id_ed25519.pub
\end{lstlisting}

  \item On GitHub: \emph{profile photo $\to$ Settings $\to$ SSH and GPG keys
  $\to$ New SSH key}. Give it a title identifying the machine (e.g.\
  ``RoboMaster Pi''), paste the full line printed above, and save.

  \item Verify the key works before attempting the clone:

\begin{lstlisting}[style=shell]
ssh -T git@github.com
\end{lstlisting}

  A successful setup answers with
  \texttt{Hi <username>! You've successfully authenticated, but GitHub does
  not provide shell access.} --- that message is the expected result, not an
  error.
\end{enumerate}

If you prefer HTTPS to SSH, use the second clone command instead and
authenticate with your GitHub username and a personal access token
\emph{that you generate for your own account} --- created under
\emph{Settings $\to$ Developer settings $\to$ Personal access tokens}, scoped
to repository read access, and never shared with anyone else or written into
a document.

\section{Building on Ubuntu 20.04 (the Raspberry Pi --- normal case)}
\label{sec:build2004}

This is the normal path: the Raspberry Pi runs Ubuntu 20.04, ROS Noetic's
officially supported target, so the build is standard and needs no
workarounds. (A lab PC also on 20.04 builds identically.)

\begin{lstlisting}[style=shell]
cd ~/carolus_ws
catkin_make --pkg ff_msgs   # build this dependency first, avoids ordering issues
catkin_make
source devel/setup.bash
\end{lstlisting}

\textbf{That is the whole build.} Skip the 22.04 section entirely --- it
exists only for a lab PC on an unsupported OS --- and go to
Section~\ref{sec:launching}.

On a Pi, expect the compile to take several minutes. If memory is tight on a
4\,GB model, pass \texttt{-j2} to limit parallel jobs.

\section{Building on Ubuntu 22.04 (lab-PC fallback only)}
\label{sec:build2204}

\textbf{Skip this section unless you are building on a lab PC running
Ubuntu 22.04.} None of it applies to the Raspberry Pi.

ROS Noetic officially targets Ubuntu~20.04 only. On 22.04, four system-level
fixes are needed before \texttt{catkin\_make} succeeds, and a fifth at launch
time. They are recorded because this project's lab PC was upgraded to 22.04
unintentionally and the workarounds took real effort to find --- not because
this is a recommended configuration.

\subsection{Fix 1: glibc symlinks}

\texttt{librt}, \texttt{libpthread}, and \texttt{libdl} were merged into
\texttt{libc.so.6} starting with glibc~2.34 (shipped in 22.04); older
binaries and build systems still look for the separate files.

\begin{lstlisting}[style=shell]
sudo ln -sf /usr/lib/x86_64-linux-gnu/libc.so.6 /usr/lib/x86_64-linux-gnu/librt.so
sudo ln -sf /usr/lib/x86_64-linux-gnu/libc.so.6 /usr/lib/x86_64-linux-gnu/libpthread.so
sudo ln -sf /usr/lib/x86_64-linux-gnu/libc.so.6 /usr/lib/x86_64-linux-gnu/libdl.so
\end{lstlisting}

\subsection{Fix 2: libpython3.8 symlink}

Referenced by \texttt{rospack} (an official ROS Noetic package):

\begin{lstlisting}[style=shell]
sudo ln -sf /usr/lib/x86_64-linux-gnu/libpython3.8.so.1 /usr/lib/x86_64-linux-gnu/libpython3.8.so
\end{lstlisting}

\subsection{Fix 3: Boost 1.71 (real libraries, not symlinks)}

ROS Noetic's official binaries need Boost~1.71; 22.04 ships 1.74, which is
ABI-incompatible for the libraries ROS actually links (regex, filesystem).
Symlinking will not work here --- the real 1.71 \texttt{.so} files are
needed:

\begin{lstlisting}[style=shell]
# Download from http://old-releases.ubuntu.com/ubuntu/pool/{main,universe}/b/boost1.71/
# the following packages: libboost-{chrono,date-time,filesystem,program-options,
# regex,system,thread}1.71.0

dpkg-deb -x <package>.deb <extract_dir>
# copy each extracted .so.1.71.0 file into /usr/lib/x86_64-linux-gnu/
sudo ldconfig
\end{lstlisting}

\textbf{Pitfall:} if a symlink already exists at the destination path
(pointing at the 1.74 library), a plain \texttt{cp} writes \emph{through}
the symlink and corrupts the real system 1.74 file. Always remove the
symlink first:

\begin{lstlisting}[style=shell]
sudo rm -f /usr/lib/x86_64-linux-gnu/libboost_<name>.so.1.74.0   # if it is a symlink
# then copy the real 1.71 file to that path
\end{lstlisting}

If this happens anyway, recovery is a straightforward package reinstall:

\begin{lstlisting}[style=shell]
sudo apt install --reinstall libboost-{chrono,date-time,filesystem,program-options,regex,system,thread}1.74.0
\end{lstlisting}

\subsection{Fix 4: ICU 67}

A dependency of Boost~1.71's regex component; 22.04 ships 66/70, not 67.

\begin{lstlisting}[style=shell]
# Download libicu67 from old-releases.ubuntu.com
# copy libicuuc.so.67*, libicui18n.so.67*, libicudata.so.67* into /usr/lib/x86_64-linux-gnu/
sudo ldconfig
\end{lstlisting}

\subsection{Building, after the four fixes}

\begin{lstlisting}[style=shell]
cd ~/carolus_ws
rm -rf build/ devel/
catkin_make --pkg ff_msgs
catkin_make
source devel/setup.bash
\end{lstlisting}

\subsection{Fix 5: LD\_PRELOAD at launch (a runtime fix, not a build fix)}

A binary compiled on 22.04 links against \texttt{liblog4cxx.so.12} and
\texttt{libtinyxml2.so.9}; ROS Noetic's own libraries
(\texttt{librosconsole.so}, \texttt{libimage\_transport.so}) expect
\texttt{.so.10}/\texttt{.so.6}. Both versions load at runtime and the 22.04
version wins global symbol resolution, which segfaults inside
\texttt{ros::console::initialize()}. The fix, scoped to the affected node in
its launch file (already applied in the launch file below, and in
\texttt{ekf\_lite.launch}):

\begin{lstlisting}[style=code,language=XML]
<env name="LD_PRELOAD"
     value="/usr/lib/x86_64-linux-gnu/liblog4cxx.so.10:/usr/lib/x86_64-linux-gnu/libtinyxml2.so.6" />
\end{lstlisting}

Apply the same pattern to any other C++ node recompiled under 22.04 that
segfaults the same way on startup.

\subsection{If an additional ROS package needs to be built from source}

Packages unavailable via \texttt{apt} on 22.04 (e.g.\ \texttt{robot\_localization},
for planned EKF sensor-fusion work) can fail with CMake errors
about a missing \texttt{<pkg>Config.cmake} or a nonexistent
\texttt{orocos\_kdl} include path --- seven packages
(\texttt{eigen\_conversions}, \texttt{tf2\_geometry\_msgs},
\texttt{kdl\_parser}, \texttt{kdl\_conversions},
\texttt{robot\_state\_publisher}, \texttt{tf2\_kdl},
\texttt{tf\_conversions}) are broken this way by the 20.04$\to$22.04
packaging difference. Without root access to fix the real files under
\texttt{/opt/ros/noetic/}, the working solution is a set of locally
corrected copies in \texttt{carolus\_ws/cmake\_shims/<pkg>\_fix/}, given
priority via \texttt{list(PREPEND CMAKE\_PREFIX\_PATH ...)} in the
\texttt{CMakeLists.txt} of whichever package is being built.

\textbf{Stale-cache pitfall:} if a shim is created only after a first failed
build attempt, \texttt{<pkg>\_DIR} may stay cached with the old broken path.
Clear just that variable rather than wiping the whole build:

\begin{lstlisting}[style=shell]
cmake -U <pkg>_DIR
\end{lstlisting}

\section{Launching Carolus}
\label{sec:launching}

Carolus's parameters are split in two, so the node can run on other robots
without editing any code:

\begin{itemize}
  \item \texttt{src/carolus\_node/config/robomaster\_s1.yaml} --- everything
  specific to \emph{this} robot and \emph{this} beacon: camera intrinsics,
  distortion, beacon geometry, blob-detection tuning.
  \item \texttt{src/carolus\_node/launch/carolus.launch} --- the generic
  launch file. It contains no hardware values at all; it loads a profile and
  wires up topics.
\end{itemize}

\texttt{testcarolus.launch} is the RoboMaster S1 entry point: it selects the
S1 profile, adds an S1-specific static transform, and delegates to
\texttt{carolus.launch}.

\subsection{The hardware profile}

The values below are this robot's. \textbf{Every one of them must be
re-measured for different hardware} --- see Chapter~\ref{ch:calibration} for
the camera intrinsics, and the note on \texttt{known\_points} at the end of
this chapter for the beacon.

\begin{lstlisting}[style=code]
# --- Camera intrinsics (MEASURE THESE FOR YOUR OWN CAMERA) ---
fx: 546.1957
fy: 547.0838
cx: 575.6041
cy: 372.1876
distortion: "-0.1479, 0.1452, -0.0023, 0.0025, 0.0070"   # k1 k2 p1 p2 k3

# --- Beacon geometry (MEASURE THIS FOR YOUR OWN BEACON) ---
# Four LED positions in metres, Carolus's left-handed frame. One point is
# deliberately off the plane of the other three -- P4P needs that.
known_points: "0.0825 0.0 0.0 -0.0825 0.0 0.0 0.0 0.072 0.0 0.0 0.0 0.0555"

# --- Blob detection ---
min_circularity: 0.6
saturation_threshold: 80
min_area: 8.0
max_area: 1800.0
max_distance_lim: 1000.0
lb_hue: 90.0        # OpenCV 0-180 scale, NOT 0-360. Blue LEDs here.
ub_hue: 140.0

# --- Preprocessing ---
kernel_size_gaussian: 3
kernel_size_morph: 3
image_threshold: 190

# --- Outlier rejection (moving window over successive poses) ---
filter_size: 7
translation_threshold: 0.5
rotation_threshold: 0.3
fifo_on: true

# --- Execution ---
num_threads: 1
queue_size: 10
benchtest: true

# --- Camera selection (Astrobee-era flags, kept for compatibility) ---
nav_cam: true
dock_cam: false
fisheye: false
fov: true
mono: false
bot_name: "wannabee"   # inert legacy label, no functional effect
\end{lstlisting}

\subsection{Launching}

\textbf{On the Raspberry Pi (the normal case).} The
\texttt{ubuntu2204\_preload:=false} argument is required: the
\texttt{LD\_PRELOAD} workaround it disables exists only for the 22.04
library mismatch (Fix 5), and it hardcodes \texttt{x86\_64} paths, so
leaving it on would be not merely unnecessary but wrong on \texttt{aarch64}.

\begin{lstlisting}[style=shell]
# On the Pi, with the camera bridge already running (Chapter 6):
source /opt/ros/noetic/setup.bash
source ~/carolus_ws/devel/setup.bash
roslaunch carolus_node testcarolus.launch ubuntu2204_preload:=false
\end{lstlisting}

Carolus and the camera bridge then both run on the Pi, and no image crosses
the network --- which is what makes the robot usable away from the lab PC,
and what gives the camera input rate measured in
Section~\ref{sec:pi-benchmark}.

\textbf{On a lab PC running Ubuntu 22.04} (the development fallback), omit
the argument so the preload stays enabled:

\begin{lstlisting}[style=shell]
roslaunch carolus_node testcarolus.launch
\end{lstlisting}

\textbf{On a different robot entirely}, copy the profile, substitute your own
measured values, and point Carolus at your own camera topic. No source file
and no launch file needs editing:

\begin{lstlisting}[style=shell]
roslaunch carolus_node carolus.launch \
    config:=$(rospack find carolus_node)/config/my_robot.yaml \
    camera_topic:=/my/camera/image_raw \
    ubuntu2204_preload:=false
\end{lstlisting}

The \texttt{robomaster\_cam} package (Chapter~\ref{ch:camera}) is the DJI SDK
bridge for this specific robot and is \textbf{not} needed on other hardware
--- supply the camera stream however your own platform already does.

Launch it (after sourcing the workspace):

\begin{lstlisting}[style=shell]
roslaunch carolus_node testcarolus.launch
\end{lstlisting}

\section{Verifying it is running}

\begin{lstlisting}[style=shell]
rosnode info /carolus_astrobee_rex
rostopic hz /pose
rostopic hz /postprocessed/image
\end{lstlisting}

With a beacon visible to the camera, \texttt{/pose} publishes a
\texttt{geometry\_msgs/PoseStamped}. Visualize blob detection in RViz by
adding an \textbf{Image} display on \texttt{/postprocessed/image}: correctly
detected blobs appear as small circles over the live feed.

\section{Warnings you can safely ignore}

Three messages appear routinely on a healthy system. They are listed here so
that nobody spends an afternoon diagnosing them, which has happened on this
project more than once.

\begin{itemize}
  \item \texttt{bash: .../catkin\_ws/devel/setup.bash: No such file or
  directory} --- a leftover line in a shell profile referring to a workspace
  that no longer exists. Harmless; remove the line from \texttt{\textasciitilde/.bashrc}
  if it bothers you.

  \item \texttt{WARNING: disk usage in log directory [...] is over 1GB} ---
  accumulated ROS logs. Clear them when convenient:

\begin{lstlisting}[style=shell]
rosclean purge
\end{lstlisting}

  \item \texttt{imageio 2.35.1 has requirement pillow>=8.3.2, but you'll have
  pillow 7.0.0} --- printed while installing the \texttt{robomaster} SDK. It
  comes from \texttt{myqr}, a transitive dependency used only for generating
  Wi-Fi QR codes, which this RNDIS-based setup never uses. The version
  conflict has no effect here.
\end{itemize}

\section{Beacon geometry and blob-filtering parameters}

The parameters most likely to need tuning for a new beacon or lighting
condition:

\begin{center}
\begin{tabular}{ll}
\toprule
Parameter & Meaning \\
\midrule
\texttt{min\_circularity} & Rejects non-round contours; closer to 1 is stricter. \\
\texttt{saturation\_threshold} & Rejects low-saturation pixels (in HSV) before hue filtering. \\
\texttt{min\_area}, \texttt{max\_area} & Rejects contours outside this pixel-area range. \\
\texttt{max\_distance\_lim} & Rejects candidate 4-blob groups whose pairwise distances exceed this. \\
\texttt{lb\_hue}, \texttt{ub\_hue} & Hue range the beacon's LEDs must fall within, in \textbf{OpenCV's 0--180 scale}, not 0--360. \\
\bottomrule
\end{tabular}
\end{center}

\textbf{The hue scale is worth stating explicitly, because getting it wrong
silently kills all detection.} OpenCV stores hue on 0--180 (half the usual
0--360 degree circle, so it fits in one byte), and the node compares the
blob's mean hue against \texttt{lb\_hue}/\texttt{ub\_hue} on that same
0--180 scale. Rough reference points: red $\approx$ 0--10 and 170--180,
green $\approx$ 50--80, blue $\approx$ 100--140. This project's beacon uses
blue LEDs, hence the 90--140 window in the launch file. If you enter a value
from the 0--360 convention (e.g.\ 240 for blue), every blob is rejected and
the node runs normally while never publishing a pose.

\texttt{known\_points} is the beacon's four LED positions in metres, in
Carolus's own left-handed frame, with one point deliberately off the plane
of the other three (needed for P4P to resolve pose unambiguously).
\textbf{These are a physical measurement of this project's specific beacon,
not a universal constant} --- the same caveat as the camera intrinsics in
Chapter~\ref{ch:calibration}. A beacon built to different dimensions needs
these four points re-measured and substituted; reusing this project's values
against a differently-sized beacon will make the P4P solve fail silently or
produce a plausible-looking but wrong pose.

\end{document}
